On s'intéresse à l'ensemble des spectateurs d'un cinéma.

On désigne par le terme \emph{JEUNE} un spectateur âgé de 18 ans ou moins.

On désigne par le terme \emph{ABONNÉ} un spectateur ayant souscrit un abonnement à ce cinéma.

Une étude a permis d'établir les résultats suivants :

\begin{itemize}[label*=\textbullet~]
	\item 60\,\% des spectateurs sont abonnés ;
	\item parmi les spectateurs abonnés, 40\,\% sont jeunes ;
	\item parmi les spectateurs n'étant pas abonnés, 20\,\% sont jeunes.
\end{itemize}

On interroge au hasard un spectateur et on considère les événements suivants :
\begin{itemize}[label*=\textbullet~]
	\item $A$ : \og \emph{le spectateur est abonné} \fg{} ;
	\item $J$ : \og \emph{le spectateur est jeune} \fg.
\end{itemize}

\begin{enumerate}
	\item Reproduire et compléter l'arbre pondéré ci-dessous décrivant la situation.

	
	\begin{center}
		\begin{tikzpicture}[yscale=-1,
			grow = right,
			level distance=2.5cm,
			level 1/.style={sibling distance=1.6cm},
			level 2/.style={sibling distance=0.8cm},
			edge from parent/.style={draw},
			edge from parent path=
			{(\tikzparentnode.east)--(\tikzchildnode.west)},
			every node/.style={fill=white}]
		\node {}
		child {
				node {$A$}
					child { node {$J$} edge from parent node {\dots} }
					child { node {$\overline{J}$} edge from parent node {\dots} }
				edge from parent node {\dots}
				}
		child {
				node {$\overline{A}$}
					child { node {$J$} edge from parent node {\dots} }
					child { node {$\overline{J}$} edge from parent node {\dots} }
				edge from parent node {\dots}
				};
		\end{tikzpicture}













	\end{center}

	\item Calculer la probabilité $P(A \cap J)$ et interpréter le résultat dans le contexte de l'exercice.

	\item Démontrer que $P(J) = 0,32$.

	\item On considère l'affirmation suivante :

	\begin{quote}
		« \textit{Si on interroge un spectateur jeune au hasard, il y a plus d'une chance sur deux qu'il soit abonné.} »
	\end{quote}

	Cette affirmation est-elle exacte ? Justifier.

	\item Le prix d'une place de cinéma est fixé selon les règles suivantes :
	\begin{itemize}[label*=\textbullet~]
		\item un spectateur non abonné paie sa place 5 euros, quel que soit son âge ;
		\item un spectateur abonné et n'étant pas jeune paie sa place 2 euros ;
		\item un spectateur abonné et jeune ne paie pas sa place : c'est gratuit.
	\end{itemize}

	On note $X$ la variable aléatoire qui, à un spectateur choisi au hasard, donne le prix qu'il doit payer pour sa place de cinéma.

	\begin{enumerate}
		\item Déterminer la loi de probabilité de la variable aléatoire $X$.
		\item Calculer l'espérance de la variable aléatoire $X$ et interpréter le résultat dans le contexte de l'exercice.
	\end{enumerate}
\end{enumerate}
